% 有理数中不收敛的柯西列：sqrt2 的十进制截断 1, 1.4, 1.41, 1.414, 1.4142；sqrt2≈1.41421 不在 Q 中。上轴 [0.9,1.6]，下轴放大 [1.41,1.42]。
\begin{tikzpicture}[
  q/.style={circle,fill=FigureBlue,inner sep=0pt,minimum size=4.2pt},
  hole/.style={circle,draw=FigureRed,fill=white,line width=.9pt,inner sep=0pt,minimum size=5.5pt},
  note/.style={font=\FigureNoteFont\fontsize{7.5}{9}\selectfont,text=FigureMuted,align=center}]
  % upper axis: x = 10*(t-0.9) cm => [0.9,1.6] -> [0,7]
  \draw[->,FigureStroke,line width=.6pt] (-.2,0)--(7.3,0);
  \foreach \t in {1.0,1.2,1.4,1.6}{\pgfmathsetmacro{\xx}{10*(\t-0.9)}\draw[FigureStroke,line width=.5pt] (\xx,-.06)--(\xx,.06); \node[below=2pt,font=\footnotesize] at (\xx,-.04) {$\t$};}
  \node[q,label={[font=\footnotesize]above:$q_0$}] at (1,0){};
  \node[q,label={[font=\footnotesize]above:$q_1$}] at (5,0){};
  \node[q] at (5.1,0){};
  \node[hole] at (5.1421,0){};
  \node[font=\footnotesize,anchor=south west] at (5.25,.08) {$\sqrt2\notin\mathbb Q$};
  % zoom guides
  \draw[FigureGrid,line width=.6pt] (5.1,-.45)--(.5,-1.25);
  \draw[FigureGrid,line width=.6pt] (5.2,-.45)--(6.5,-1.25);
  \draw[FigureGrid,line width=.6pt] (5.1,-.12) rectangle (5.2,.12);
  % lower axis: [1.41,1.42] -> [0.5,6.5], x = 0.5 + 600*(t-1.41)
  \begin{scope}[yshift=-1.55cm]
  \draw[FigureStroke,line width=.6pt] (.3,0)--(6.7,0);
  \foreach \t/\l in {1.41/1.41,1.415/1.415,1.42/1.42}{\pgfmathsetmacro{\xx}{0.5+600*(\t-1.41)}\draw[FigureStroke,line width=.5pt] (\xx,-.06)--(\xx,.06); \node[below=2pt,font=\footnotesize] at (\xx,-.04) {$\l$};}
  \node[q,label={[font=\footnotesize]above:$q_2$}] at (0.5,0){};
  \node[q,label={[font=\footnotesize]above:$q_3$}] at (2.9,0){};
  \node[q] at (3.02,0){};
  \node[hole] at (3.028,0){};
  \node[font=\footnotesize,anchor=south west] at (3.12,.1) {$q_4,q_5,\dots$};
  \node[note] at (3.5,-.75) {后面的项越挤越近，却在$\mathbb Q$中无处落脚};
  \end{scope}
\end{tikzpicture}
